\begin{statement}{E17}{Refuted}{The embedded-component omission in the divisor-twisting lemma}
CDP98, Proposition~1.1, in its printed generality of arbitrary
nonreduced compact spaces, is false. The exact replacement is as
follows. For an effective Cartier divisor $D$ on a compact complex
space $Z$ and a locally free sheaf $\mathcal V$, the spaces
$H^0(Z,\mathcal V(-kD))$, viewed inside $H^0(Z,\mathcal V)$,
stabilize for $k\gg0$ at
\[
 \{s\in H^0(Z,\mathcal V):\operatorname{supp}(s)\cap|D|
       =\varnothing\}.
\]
In particular they eventually vanish when $Z$ has no embedded
associated components and $D$ meets every irreducible component.
This last case applies to every full Cartier fibre used in E2.
\end{statement}
\begin{proof}
Let $C=\mathbb P^1$, let $p,q$ be distinct points, and form the
square-zero complex space
\[
 |Z|=C,\qquad\mathcal O_Z=\mathcal O_C\oplus\mathbb C_p,
 \qquad(a,u)(b,v)=(ab,av+bu).
\]
Near $p$ its analytic local ring is
$\mathbb C\{z,e\}/(ze,e^2)$, and away from $p$ it is the
ordinary smooth curve. These charts glue through the quotient
by $e$, so this specifies a compact connected complex space.
It has one irreducible component, of positive dimension, and an
embedded associated point at $p$. The divisor $D=q$ is effective
Cartier: its local equation is a coordinate at $q$ and a unit
near $p$. For every integer $k$ there is the exact, split
description as an $\mathcal O_C$-module
\[
 \mathcal O_Z(-kD)=\mathcal O_C(-kq)\oplus\mathbb C_p.
\]
It gives every cohomology group:
\[
 h^0=\max(1-k,0)+1,\qquad h^1=\max(k-1,0),\qquad
 h^j=0\quad(j\geq2).
\]
In particular the nonzero section of the square-zero summand
survives for every $k\geq1$, refuting the printed proposition.

For the replacement, the descending spaces in the statement are
finite dimensional, hence stabilize. A section in their
intersection belongs to $\bigcap_k I_D^k\mathcal V_z=0$ at
every $z\in|D|$, by Krull intersection in the analytic local
ring. Its coherent support is therefore disjoint from $D$.
Conversely a section with support disjoint from $D$ vanishes
near $D$ and is a section of $\mathcal V(-kD)$ for every $k$.
This proves the exact stable-space formula. When there are no
embedded associated components, every nonzero submodule of a
locally free sheaf has support containing an irreducible
component of $Z$. Such a support cannot avoid $D$ under the
stated intersection hypothesis, proving the vanishing.
\end{proof}

\paragraph{Audit of the printed non-holomorphic proof.}
We now give its entire chain of reductions and identify its
dependencies. CDP20, Section~3, first uses constancy of the index
under Picard twists and Riemann--Roch to deduce its parts (b)
and (c) from generic vanishing. Serre duality reduces generic
vanishing to degrees zero and two. Degree zero invokes CDP98,
Corollary~1.3; that corollary uses only Proposition~1.1 on the
smooth manifold, a multiple of an effective divisor lying in
$\operatorname{Pic}^0$, and semicontinuity. E17 proves the
version of Proposition~1.1 needed there, and E2 spells out the
exponential-family semicontinuity argument. Thus this dependency
is valid and involves none of the cotangent or monodromy steps.

For degree two the printed proof resolves the meromorphic pencil
by $\sigma:\widehat Y\to Y$, with a connected map
$h:\widehat Y\to\mathbb P^1$, and writes
\[
 h^*\mathcal O(1)=\sigma^*P\otimes\mathcal O(-E),
\]
where $E$ is the exceptional divisor of the pulled-back base
ideal. The map on second cohomology from $Y$ to $\widehat Y$
is an isomorphism, since a composition of smooth blowups has
$R\sigma_*\mathcal O=\mathcal O_Y$. The divisor sequence
reduces vanishing on $\widehat Y$ to vanishing after $-tE$
and to $H^2(tE,\sigma^*W)=0$. The printed proof asserts
the exceptional higher-image vanishing needed for this step.
For this particular pencil divisor it has the following direct
justification. Pull back the evaluation sequence on $\mathbb P^1$
and tensor by $\sigma^*P^{-t}$ to obtain, for $t\geq1$,
\[
\begin{split}
0\longrightarrow&\bigl(\mathcal O(E)\otimes
              \sigma^*P^{-(t+1)}\bigr)^{\oplus t}
 \longrightarrow(\sigma^*P^{-t})^{\oplus(t+1)}\\
 &\longrightarrow\mathcal O(-tE)\longrightarrow0.
\end{split}
\]
The evaluation sequence remains exact under pullback because it
is a sequence of vector bundles with locally free quotient.
The middle term has zero $R^2\sigma_*$, by the projection
formula and $R\sigma_*\mathcal O=\mathcal O_Y$.
Every coherent sheaf has zero $R^3\sigma_*$ because the
fibres have dimension at most two. Hence
$R^2\sigma_*\mathcal O(-tE)=0$; the divisor sequence now gives
\[
 R^1\sigma_*\mathcal O_{tE}=
 R^2\sigma_*\mathcal O_{tE}=0.
\]
The support of $\sigma_*\mathcal O_{tE}$ has dimension at most
one, so Leray and the projection formula give the required
$H^2(tE,\sigma^*W)=0$. No identification of this pushforward
with a structure sheaf is needed.

This special verification matters. The analogous assertion for
an arbitrary exceptional divisor is false. Blow up a point
$p\in\mathbb P^3$ and then a smooth plane cubic in the exceptional
$\mathbb P^2$. For the second exceptional divisor $S$, projection
to that cubic is a $\mathbb P^1$-bundle; thus
$h^\bullet(S,\mathcal O_S)=(1,1,0)$ and the composite contraction
to $\mathbb P^3$ has
$R^1\sigma_*\mathcal O_S=\mathbb C_p$.
This example, also given in \cite{LRS2019}, concerns arbitrary
exceptional divisors and does not refute the just-proved assertion
for the divisor of a resolved pencil. The graph proof of E2
avoids the distinction by computing its own higher direct images
from its hypersurface equation.

The printed argument then replaces its twist $M$ by
$M\otimes P^{t+k}$ and applies Leray to $h$. It must prove
\[
 R^2h_*\sigma^*(W\otimes M\otimes P^k)=0
\]
and then the vanishing of $H^1$ of the corresponding $R^1$
twisted by $\mathcal O(t)$; the latter is ordinary Serre
vanishing on $\mathbb P^1$. For the former it uses top-degree
base change and fibrewise Serre duality. It invokes CDP98,
Proposition~1.1, first on components meeting $E$, and then
successively on components at increasing graph distance from
them. Coefficients in an auxiliary effective divisor are chosen
in reverse order of this distance, and the power of $E$ is
chosen last. E2 gives this argument with the actual pullback
maps, full multiplicity filtrations, and a finite-support
argument establishing uniformity. These fibres are Cartier in
a smooth threefold and have no embedded components, so E17's
proved version of the twisting lemma applies.

These are the only earlier CDP results invoked in Section~3:
CDP98, Corollary~1.3 and Proposition~1.1. Its opening reference
to CDP98, Theorem~2.1 identifies the proof being rewritten,
rather than using that theorem as a premise. CDP98, Lemma~1.5
is not invoked; its erroneous replacement of fibrewise vanishing
by generic vanishing is precisely what the full component
argument removes. No result from CDP20, Sections~4--7 is used.
In particular the cotangent-filtration and monodromy failures
treated in Parts A and D do not enter E2. The supplied LRS
graph proof has the same dependency on the no-embedded-component
twisting lemma, and supplies its own higher-image calculation.
Its displayed adjunction formula has a sign misprint:
for a Cartier fibre $F\subset Y$ it is
$\omega_F=(K_Y+F)|_F$, not $(K_Y-F)|_F$.
E2 uses the actual dualizing line and includes nonreduced fibres,
so that misprint has no role in the proof here.
